\documentclass[10pt,a4paper]{article}
\usepackage[utf8]{inputenc}
\usepackage[T1]{fontenc}
\usepackage[french]{babel}
\usepackage{lmodern}
\usepackage[margin=2.1cm]{geometry}
\usepackage{microtype}
\usepackage{amsmath}
\usepackage{booktabs,tabularx,array}
\usepackage{xcolor}
\usepackage{fancyhdr}
\usepackage[hidelinks]{hyperref}
\definecolor{accent}{RGB}{34,65,91}
\newcolumntype{Y}{>{\raggedright\arraybackslash}X}
\setlength{\parindent}{0pt}
\setlength{\parskip}{0.4em}
\setlength{\tabcolsep}{5pt}
\renewcommand{\arraystretch}{1.16}
\pagestyle{fancy}
\fancyhf{}
\fancyhead[L]{\small Identification des utilisateurs de Copilote}
\fancyhead[R]{\small BE MOD 7.2}
\fancyfoot[L]{\small Version de travail -- 12 octobre 2026}
\fancyfoot[R]{\thepage}
\setlength{\headheight}{14pt}
\newcommand{\afaire}[1]{\par\smallskip\textcolor{accent}{\textbf{À réaliser.}} #1\par}
\newcommand{\sourceexploration}{\textit{Source : sorties enregistrées de \texttt{datamanip/001.ipynb}, jeu d'entraînement.}}
\hypersetup{pdftitle={Identification des utilisateurs de Copilote -- Rapport de BE},pdflang={fr-FR}}

\begin{document}
\thispagestyle{plain}
\begin{center}
{\Large\bfseries\color{accent} Identification des utilisateurs de Copilote}\par
\vspace{0.35em}
{\large Rapport de bureau d'étude}\par
MOD 7.2 -- Introduction à la science des données\\
Compétition Kaggle -- Groupe 2\\[0.5em]
\textit{Version de travail du 12 octobre 2026}
\end{center}

\textbf{État du travail.} Cette première version présente le problème, les données
et l'exploration initiale déjà effectuée. Aucun modèle n'a été entraîné et aucun
score de classification n'est disponible. Les sections consacrées à la suite
de l'étude indiquent explicitement les travaux restant à mener. Les auteurs
seront renseignés lors de la finalisation du rapport.

\section{Introduction}
Copilote est un logiciel édité par Infologic pour le secteur agroalimentaire.
Son utilisation produit des traces décrivant les actions réalisées au cours
d'une session, ainsi que certains éléments de leur contexte. Ces traces
permettent d'étudier les habitudes d'utilisation du logiciel : fonctionnalités
sollicitées, écrans consultés et interactions effectuées.

Le projet pose la question suivante : peut-on reconnaître un utilisateur à
partir d'une session d'utilisation ? Dans le contexte présenté par le BE,
les testeurs internes utilisent différents profils pour tester le logiciel.
Identifier le testeur à partir de ses traces pourrait faciliter le signalement
des problèmes. La reconnaissance des habitudes pourrait aussi contribuer
à repérer des comportements inhabituels sur un compte. Cette dernière
application constitue une motivation ; elle n'est pas évaluée directement
dans les données disponibles.

\section{Problématique et objectifs}
\subsection{Une classification supervisée multiclasse}
L'unité d'observation est la \textbf{session}. La cible est l'identifiant de
l'utilisateur, nommé \texttt{util} dans le tableau d'exploration. Le jeu
d'entraînement comporte 247 utilisateurs distincts : il s'agit donc d'une
classification supervisée à 247 classes. À partir du navigateur et de la
séquence de traces d'une session, la solution devra prédire un identifiant
utilisateur d'origine.

L'enjeu est de construire une représentation qui distingue les habitudes
des utilisateurs malgré la variabilité de leurs sessions. Une fréquence
élevée dans l'ensemble du jeu ne suffit pas à prouver qu'une action est
discriminante. Il faudra mesurer l'apport des caractéristiques sur des
sessions réservées à la validation.

\subsection{Cadre du rendu}
Les consignes demandent un rapport PDF de 5 à 8 pages, un code Python
réexécutable avec ses dépendances et les prédictions sur le jeu de test.
Le dépôt des prédictions est prévu sur Moodle et sur la compétition Kaggle
du groupe 2. L'échéance annoncée dans le notebook de consignes est le
\textbf{9 novembre 2026 au soir}. Le rapport final devra justifier les
caractéristiques et les méthodes, préciser les métriques et discuter les
résultats et leurs limites.

\newpage
\section{Présentation des données et préparation initiale}
\subsection{Organisation des fichiers}
Les données sont fournies sous forme de fichiers texte séparés par des
virgules, sans en-tête. Chaque ligne représente une session ; le nombre
de champs varie selon la longueur de la séquence. Dans l'entraînement,
les deux premiers champs sont l'identifiant utilisateur et le navigateur.
Dans le test, l'identifiant est absent et le navigateur est le premier champ.

\begin{center}
\begin{tabularx}{\linewidth}{l r Y}
\toprule
Fichier & Sessions & Rôle \\
\midrule
\texttt{train.csv} & 3\,279 & Apprentissage et validation ; identités connues. \\
\texttt{test.csv} & 324 & Prédictions à produire ; identités non fournies. \\
\texttt{sample\_submission.csv} & 324 & Exemple du format de soumission. \\
\bottomrule
\end{tabularx}
\end{center}
\textit{Source : README du projet. L'exploration décrite ici porte uniquement sur le train.}

\subsection{Contenu d'une trace}
Une séquence mêle des actions et des marqueurs temporels. Certaines actions
portent des informations supplémentaires, repérées par leurs délimiteurs.
Ces informations ne sont pas nécessairement présentes dans chaque action.

\begin{center}
\begin{tabularx}{\linewidth}{l Y}
\toprule
Élément & Signification retenue \\
\midrule
Nom de l'action & Interaction, par exemple double-clic ou saisie dans un champ. \\
\texttt{(\ldots)} & Écran ou contexte de l'action. \\
\texttt{<\ldots>} & Configuration d'écran. \\
\texttt{\$\ldots\$} & Chaîne de fiche, regroupement logique dans le logiciel. \\
Suffixe \texttt{1} & Action effectuée en modification dans la fiche. \\
\texttt{tXX} & Repère de fenêtre temporelle exprimé en secondes. \\
\bottomrule
\end{tabularx}
\end{center}

Le nombre total de traces inclut les marqueurs temporels : il ne doit donc
pas être confondu avec le nombre d'actions. Le plus grand repère temporel
est utilisé comme indicateur descriptif, sans être assimilé à une durée
exacte de session. Des repères peuvent être répétés.

\subsection{Conservation et contrôles des données}
Le notebook lit les lignes une à une et conserve les traces sous forme de
listes. Il construit \texttt{df\_raw} avec quatre colonnes :
\texttt{session\_id}, \texttt{util}, \texttt{navigateur} et \texttt{traces}.
Le numéro de ligne sert de repère et ne constitue pas une caractéristique
d'identification prévue. Les transformations sont effectuées dans
\texttt{df\_train}. Les listes de traces sont également copiées pour éviter
de partager des objets modifiables avec le tableau brut.

Les contrôles enregistrés ne détectent aucune valeur manquante dans les
quatre colonnes initiales, aucun identifiant ou navigateur vide et aucun
champ de trace vide. Aucun doublon strict n'est observé en comparant
l'utilisateur, le navigateur et la séquence complète, sans le numéro de
ligne. Cela n'exclut pas des sessions proches ni des anomalies de contenu.
Les vérifications finales confirment le même nombre de sessions et la
conservation des colonnes brutes dans la copie de travail.

\sourceexploration

\newpage
\section{Exploration du jeu d'entraînement}
\subsection{Répartition des utilisateurs et navigateurs}
Les 3\,279 sessions se répartissent entre 247 classes. Le nombre de sessions
par utilisateur varie de 4 à 75, avec une médiane de 12 et une moyenne
de 13,28. Les premier et troisième quartiles sont respectivement 7 et 16.
Le déséquilibre des classes et le faible nombre d'exemples de certains
utilisateurs devront être pris en compte dans la validation.

\begin{center}
\begin{tabular}{lrr}
\toprule
Navigateur & Sessions & Part des sessions \\
\midrule
Firefox & 1\,466 & 44,71\,\% \\
Google Chrome & 1\,339 & 40,84\,\% \\
Microsoft Edge & 451 & 13,75\,\% \\
Opera & 23 & 0,70\,\% \\
\bottomrule
\end{tabular}
\end{center}

Firefox et Chrome regroupent 85,55\,\% des sessions. Le croisement entre
utilisateur et navigateur montre qu'un seul navigateur est observé pour
chacun des 247 utilisateurs dans le train. Plusieurs utilisateurs partagent
toutefois le même navigateur : cet indice ne suffit pas à identifier une
personne. Sa stabilité dans de nouvelles sessions n'est pas démontrée.

\subsection{Longueur des sessions}
Les marqueurs reconnus par l'expression régulière \texttt{t\textbackslash d+}
sur la totalité de la trace sont séparés des actions. Le notebook calcule
ensuite le nombre d'actions, le nombre de marqueurs et le plus grand repère.

\begin{center}
\small
\begin{tabular}{lrrrrr}
\toprule
Indicateur & Minimum & Q1 & Médiane & Q3 & Maximum \\
\midrule
Actions & 0 & 153 & 340 & 818,5 & 11\,977 \\
Marqueurs temporels & 1 & 38 & 82 & 197,5 & 2\,491 \\
Dernier repère (s) & 5 & 185 & 395 & 945 & 11\,790 \\
\bottomrule
\end{tabular}
\end{center}

La moyenne du nombre d'actions est de 692,36, contre une médiane de 340.
Cette différence et le maximum de 11\,977 actions indiquent une forte
asymétrie des longueurs. Des comptages bruts risquent ainsi de refléter
la longueur de la session autant que les habitudes de son utilisateur.
La comparaison avec des fréquences normalisées est une piste pour la suite.

\subsection{Session sans action}
Une session, repérée par \texttt{session\_id = 1252}, appartient à
l'utilisateur \texttt{fvc}, sous Google Chrome, et contient uniquement
\texttt{t5}, \texttt{t10} et \texttt{t15}. Elle est conservée à ce stade.
L'absence d'action n'est pas un champ vide : la session contient des
informations temporelles et un navigateur, mais aucun contenu d'action.

La future chaîne de traitement devra accepter ce cas et gérer les ratios
dont le dénominateur est nul. L'effet de cette session sur l'apprentissage
ou l'évaluation n'a pas encore été mesuré. Sa conservation évite une
suppression décidée sans analyse de son incidence.

\sourceexploration

\newpage
\subsection{Types d'actions et contextes}
Le type d'action est obtenu en séparant le nom des informations de contexte
au premier délimiteur parenthèse, chevron ou dollar, puis en retirant
l'éventuel suffixe final \texttt{1}. Les contextes sont extraits à l'aide
des délimiteurs présents dans la trace, sans compléter les valeurs absentes.

\begin{center}
\begin{tabularx}{\linewidth}{Y r Y}
\toprule
Information & Valeurs distinctes & Valeur la plus fréquente \\
\midrule
Types d'actions & 35 & Exécution d'un bouton \\
Écrans ou contextes & 418 & \texttt{MAINT} \\
Configurations & 219 & \texttt{DEF\_03/24} \\
Chaînes de fiches & 76 & \texttt{VT} \\
\bottomrule
\end{tabularx}
\end{center}

L'exécution d'un bouton représente 510\,645 occurrences et la saisie dans
un champ 265\,633 occurrences dans l'ensemble du train. Ces comptages
portent sur des occurrences, et non sur le nombre de sessions ou
d'utilisateurs. Ils sont sensibles aux sessions longues et aux répétitions.
Ils ne constituent pas des mesures de l'importance prédictive des actions.

Les écrans, configurations et chaînes peuvent renseigner sur les activités
ou les contextes de travail propres à certains utilisateurs. Cette utilité
reste une hypothèse. La présence de nombreuses modalités suggère une
représentation potentiellement creuse, dont la construction et le traitement
des modalités nouvelles devront être définis avant l'entraînement.

\subsection{Actions en modification}
Le notebook compte comme action en modification toute action terminée
par \texttt{1}. Pour une session $s$ ayant $n_s$ actions, dont $m_s$ en
modification, la proportion est définie par :
\[
 p_s=\frac{m_s}{n_s}\quad\text{si } n_s>0.
\]
Pour la session sans action, cette proportion est laissée manquante :
elle est indéfinie et ne doit pas être interprétée comme une proportion nulle.

\begin{center}
\begin{tabularx}{\linewidth}{Y r}
\toprule
Indicateur & Valeur \\
\midrule
Médiane du nombre d'actions en modification par session & 34 \\
Moyenne du nombre d'actions en modification par session & 108,30 \\
Médiane des proportions par session & 9,09\,\% \\
Moyenne des proportions par session & 15,54\,\% \\
Sessions avec une proportion définie & 3\,278 \\
Sessions avec une proportion non définie & 1 \\
\bottomrule
\end{tabularx}
\end{center}

La moyenne de 15,54\,\% est la moyenne des proportions sur les sessions
dont le ratio est défini. Chaque session y a le même poids ; elle ne
correspond pas à la proportion obtenue en regroupant toutes les actions.
Le comportement en modification constitue une caractéristique candidate,
dont l'apport devra être évalué conjointement avec les autres variables.

\subsection{Bilan de l'exploration}
Les données sont exploitables sur le plan des contrôles de structure
effectués, mais présentent des longueurs très variables et des classes
inégalement représentées. Le navigateur, la composition des actions et
leurs contextes fournissent des pistes de caractérisation. À ce stade,
aucune observation descriptive ne permet de conclure à la fiabilité
d'un système d'identification.

\sourceexploration

\newpage
\section{Méthodologie envisagée}
\afaire{Construire une représentation par session et comparer son apport
en validation. Le tableau ci-dessous décrit des pistes, pas des choix
déjà validés ni des modèles déjà entraînés.}

\begin{center}
\small
\begin{tabularx}{\linewidth}{l Y}
\toprule
Famille & Caractéristiques candidates et question à examiner \\
\midrule
Informations globales & Navigateur, nombre d'actions, marqueurs et dernier repère ; apport du contexte et de la longueur. \\
Actions & Comptages et fréquences par type ; intérêt de normaliser les sessions longues. \\
Contextes & Présence ou fréquence des écrans, configurations et chaînes ; apport des usages métiers. \\
Modification & Nombre et proportion d'actions en modification ; traitement explicite du ratio indéfini. \\
Séquences & Motifs courts ou transitions ; apport éventuel de l'ordre des actions. \\
\bottomrule
\end{tabularx}
\end{center}

Une référence simple, par exemple la prédiction de la classe la plus
fréquente de l'apprentissage, permettra de situer les performances. Des
classifieurs adaptés à des variables numériques et catégorielles ou à
des représentations creuses seront ensuite comparés. La liste finale
des méthodes, leurs paramètres et les transformations seront renseignés
après les expériences. L'identité cible et le numéro de ligne devront
être exclus des variables explicatives.

\section{Validation et métriques}
\afaire{Fixer une séparation stratifiée des sessions étiquetées, vérifier
la représentation de toutes les classes et consigner la graine aléatoire.
Avec au moins quatre sessions par classe, une validation croisée à quatre
plis est une option à examiner ; les scores des classes rares resteront
sensibles au petit nombre d'exemples.}

Les vocabulaires, encodages, normalisations et éventuelles imputations
devront être ajustés sur la partie apprentissage de chaque séparation.
Les données de validation serviront à comparer les solutions ; le test
sans étiquettes ne permet pas de calculer un score local. Une séparation
aléatoire par session évalue l'identification sur des sessions réservées,
sans démontrer une robustesse à un changement temporel ou d'environnement.

L'accuracy mesurera la proportion de sessions correctement classées.
Le F1 combine précision et rappel, avec $F_1=2PR/(P+R)$ lorsque le
dénominateur est non nul. Le README mentionne un F1 moyen pour Kaggle,
mais son mode de moyenne reste à confirmer. Un F1 macro, moyenne non
pondérée des F1 par classe, est envisagé en complément pour examiner
le déséquilibre ; il ne sera pas présenté comme la métrique officielle
tant que ce point n'aura pas été vérifié.

\section{Résultats de classification}
\afaire{Présenter les scores de référence et des modèles sur le même
protocole, les paramètres retenus et la variabilité éventuelle entre plis.
Aucun score, classement Kaggle ou fichier de prédictions n'est disponible
à ce stade. Les tableaux descriptifs précédents sont les seuls résultats
actuellement documentés.}

\section{Discussion et conclusion provisoire}
\afaire{Analyser les erreurs, les classes rares et l'apport des familles
de variables par des comparaisons avec et sans ces variables. Discuter
la dépendance au navigateur et aux contextes, ainsi que la session sans
action. Documenter l'environnement effectivement testé : les consignes
évoquent Python 3.14, tandis que le dépôt indique actuellement Python 3.11.}

L'exploration fournit une base pour construire la chaîne de classification.
La prochaine étape est de définir un protocole de validation et une première
représentation des sessions, puis de mesurer leurs performances. La
solution retenue devra enfin produire 324 prédictions, dans l'ordre du
test, aux colonnes \texttt{RowId} (de 1 à 324) et \texttt{prediction}.

\vfill
{\footnotesize\textbf{Documents de référence locaux :}
\texttt{README.md} ; \texttt{Notebooks/BE\_CS1\_instructions.ipynb} ;
\texttt{datamanip/001.ipynb}. Les constats d'exploration proviennent des
sorties enregistrées du dernier notebook ; aucune nouvelle expérience
de classification n'a été exécutée pour cette rédaction.}
\end{document}
